% GENERATED FILE. Edit canonical lessons, then run npm run generate:books.
% canonical-source-hash: fnv1a64:0d1beefb7d7bad04
% canonical-lessons: FA-C490-sarmakhordegi, FA-C490-anfolanza, FA-C490-atse-kardan, FA-C490-galudard, FA-C490-tab-kardan

\chapter{Colds and Flu}
\label{ch:fa-colds-and-flu}
\hlchaptermodality{490}

\begin{chapteropening}
\textbf{By the end of this chapter:} \emph{I can say I have a cold or the flu.}

Complete the last lesson of chapter 490: I can say I have a cold or the flu.
\end{chapteropening}

\section[\texorpdfstring{\fa{سرماخوردگی} (sarmâkhordegi)}{سرماخوردگی (sarmâkhordegi)}]{\fa{سرماخوردگی} (sarmâkhordegi) — a cold (illness)}

\label{lesson:FA-C490-sarmakhordegi}



\begin{quote}
Before the new one: say the Persian for shy, then the Persian for to make fun of.
\end{quote}

\subsection*{You'll want to know: \fa{سرماخوردگی}}
\textbf{\fa{سرماخوردگی}} (\emph{sarmâkhordegi}) — \textquotedblleft{}a cold (illness)\textquotedblright{}.

\begin{quote}
\fa{او} \fa{سرماخوردگی} \fa{دارد} \fa{و} \fa{سر} \fa{کار} \fa{نمیرود}.
\end{quote}

\begin{quote}
\emph{u sarmâkhordegi dârad va sar-e kâr nemiravad.}
\end{quote}

\begin{quote}
He has a cold and isn't going to work.
\end{quote}

The noun from \emph{sarmâ khordan} 'to catch a cold', with the ending \emph{-egi}.

\begin{cousinweb}
\emph{sarmâkhordegi} is built on native Persian pieces: \emph{sarmâ}, 'cold', \emph{khord}, the past stem of \emph{khordan}, 'to eat', and the noun ending \emph{-egi}: the state of having eaten the cold.
\end{cousinweb}

\begin{grammarlens}[title={What you've built}]
One more word for telling, explaining and describing at B1.
\end{grammarlens}

\subsection*{Your turn}
\begin{itemize}
\raggedright
  \item \emph{Say it:} \emph{sarmâkhordegi}
  \item \emph{Say it:} \emph{sarmâkhordegi}, once more
  \item \emph{Say it:} a sentence of your own with \emph{sarmâkhordegi}
  \item \emph{Recall:} say the Persian for shy, then the Persian for to make fun of, then say \emph{sarmâkhordegi} again
\end{itemize}

\begin{tcolorbox}[breakable,colback=teal!4,colframe=teal!35!black,title={Before you move on}]
What does \fa{سرماخوردگی} mean? (\textquotedblleft{}A cold (illness)\textquotedblright{}.) Say it once more.
\end{tcolorbox}
\section[\texorpdfstring{\fa{آنفولانزا} (ânfolânzâ)}{آنفولانزا (ânfolânzâ)}]{\fa{آنفولانزا} (ânfolânzâ) — flu, influenza}

\label{lesson:FA-C490-anfolanza}



\begin{quote}
Before the new one: say the Persian for to make fun of, then the Persian for a cold.
\end{quote}

\subsection*{You'll want to know: \fa{آنفولانزا}}
\textbf{\fa{آنفولانزا}} (\emph{ânfolânzâ}) — \textquotedblleft{}flu, influenza\textquotedblright{}.

\begin{quote}
\fa{این} \fa{زمستان} \fa{خیلی} \fa{از} \fa{بچه} \fa{ها} \fa{آنفولانزا} \fa{گرفتند}.
\end{quote}

\begin{quote}
\emph{in zemestân kheyli az bachchehâ ânfolânzâ gereftand.}
\end{quote}

\begin{quote}
This winter a lot of children caught flu.
\end{quote}

\emph{ânfolânzâ gereftan} — to catch flu.

\begin{cousinweb}
\emph{ânfolânzâ} is a modern loan, usually traced to French \emph{influenza}, which French took from Italian. EtymDB has no entry for the Persian word.
\end{cousinweb}

\begin{grammarlens}[title={What you've built}]
One more word for telling, explaining and describing at B1.
\end{grammarlens}

\subsection*{Your turn}
\begin{itemize}
\raggedright
  \item \emph{Say it:} \emph{ânfolânzâ}
  \item \emph{Say it:} \emph{ânfolânzâ}, once more
  \item \emph{Say it:} a sentence of your own with \emph{ânfolânzâ}
  \item \emph{Recall:} say the Persian for to make fun of, then the Persian for a cold, then say \emph{ânfolânzâ} again
\end{itemize}

\begin{tcolorbox}[breakable,colback=teal!4,colframe=teal!35!black,title={Before you move on}]
What does \fa{آنفولانزا} mean? (\textquotedblleft{}Flu, influenza\textquotedblright{}.) Say it once more.
\end{tcolorbox}
\section[\texorpdfstring{\fa{عطسه} \fa{کردن} ('atse kardan)}{عطسه کردن ('atse kardan)}]{\fa{عطسه} \fa{کردن} ('atse kardan) — to sneeze}

\label{lesson:FA-C490-atse-kardan}



\begin{quote}
Before the new one: say the Persian for a cold, then the Persian for flu.
\end{quote}

\subsection*{You'll want to know: \fa{عطسه} \fa{کردن}}
\textbf{\fa{عطسه} \fa{کردن}} (\emph{'atse kardan}) — \textquotedblleft{}to sneeze\textquotedblright{}.

\begin{quote}
\fa{او} \fa{در} \fa{بهار} \fa{زیاد} \fa{عطسه} \fa{میکند}.
\end{quote}

\begin{quote}
\emph{u dar bahâr ziyâd 'atse mikonad.}
\end{quote}

\begin{quote}
She sneezes a lot in spring.
\end{quote}

\emph{'atse} 'a sneeze' + \emph{kardan}; present stem \emph{kon-}.

\begin{cousinweb}
\emph{'atse} is a loan, usually traced to Arabic \emph{ʿaṭsa}, 'a sneeze'. EtymDB has no entry for it. \emph{kardan}, 'to do', is inherited Persian.
\end{cousinweb}

\begin{grammarlens}[title={What you've built}]
One more word for telling, explaining and describing at B1.
\end{grammarlens}

\subsection*{Your turn}
\begin{itemize}
\raggedright
  \item \emph{Say it:} \emph{'atse kardan}
  \item \emph{Say it:} \emph{'atse kardan}, once more
  \item \emph{Say it:} a sentence of your own with \emph{'atse kardan}
  \item \emph{Recall:} say the Persian for a cold, then the Persian for flu, then say \emph{'atse kardan} again
\end{itemize}

\begin{tcolorbox}[breakable,colback=teal!4,colframe=teal!35!black,title={Before you move on}]
What does \fa{عطسه} \fa{کردن} mean? (\textquotedblleft{}To sneeze\textquotedblright{}.) Say it once more.
\end{tcolorbox}
\section[\texorpdfstring{\fa{گلودرد} (galudard)}{گلودرد (galudard)}]{\fa{گلودرد} (galudard) — a sore throat}

\label{lesson:FA-C490-galudard}



\begin{quote}
Before the new one: say the Persian for flu, then the Persian for to sneeze.
\end{quote}

\subsection*{You'll want to know: \fa{گلودرد}}
\textbf{\fa{گلودرد}} (\emph{galudard}) — \textquotedblleft{}a sore throat\textquotedblright{}.

\begin{quote}
\fa{از} \fa{دیروز} \fa{گلودرد} \fa{دارم}.
\end{quote}

\begin{quote}
\emph{az diruz galudard dâram.}
\end{quote}

\begin{quote}
I've had a sore throat since yesterday.
\end{quote}

\emph{galu} 'throat' + \emph{dard} 'pain', on the pattern of \emph{sardard} 'headache'.

\begin{cousinweb}
\emph{galudard} is built on two inherited words: \emph{galu}, 'throat', usually traced to Middle Persian \emph{galōg}, and \emph{dard}, 'pain', as in \emph{sardard}.
\end{cousinweb}

\begin{grammarlens}[title={What you've built}]
One more word for telling, explaining and describing at B1.
\end{grammarlens}

\subsection*{Your turn}
\begin{itemize}
\raggedright
  \item \emph{Say it:} \emph{galudard}
  \item \emph{Say it:} \emph{galudard}, once more
  \item \emph{Say it:} a sentence of your own with \emph{galudard}
  \item \emph{Recall:} say the Persian for flu, then the Persian for to sneeze, then say \emph{galudard} again
\end{itemize}

\begin{tcolorbox}[breakable,colback=teal!4,colframe=teal!35!black,title={Before you move on}]
What does \fa{گلودرد} mean? (\textquotedblleft{}A sore throat\textquotedblright{}.) Say it once more.
\end{tcolorbox}
\section[\texorpdfstring{\fa{تب} \fa{کردن} (tab kardan)}{تب کردن (tab kardan)}]{\fa{تب} \fa{کردن} (tab kardan) — to run a fever}

\label{lesson:FA-C490-tab-kardan}



\begin{quote}
Before the new one: say the Persian for to sneeze, then the Persian for a sore throat.
\end{quote}

\subsection*{You'll want to know: \fa{تب} \fa{کردن}}
\textbf{\fa{تب} \fa{کردن}} (\emph{tab kardan}) — \textquotedblleft{}to run a fever\textquotedblright{}.

\begin{quote}
\fa{بچه} \fa{دیشب} \fa{تب} \fa{کرد}.
\end{quote}

\begin{quote}
\emph{bachche dishab tab kard.}
\end{quote}

\begin{quote}
The child ran a fever last night.
\end{quote}

\emph{tab} 'fever' + \emph{kardan}. \emph{tab dâshtan} — to have a temperature.

\begin{cousinweb}
\emph{tab} is a native Persian word for fever; EtymDB lists Armenian \emph{vap}, 'heat', beside it (EtymDB). \emph{kardan}, 'to do', is inherited Persian.
\end{cousinweb}

\begin{grammarlens}[title={What you've built}]
One more word for telling, explaining and describing at B1.
\end{grammarlens}

\subsection*{Your turn}
\begin{itemize}
\raggedright
  \item \emph{Say it:} \emph{tab kardan}
  \item \emph{Say it:} \emph{tab kardan}, once more
  \item \emph{Say it:} a sentence of your own with \emph{tab kardan}
  \item \emph{Recall:} say the Persian for to sneeze, then the Persian for a sore throat, then say \emph{tab kardan} again
\end{itemize}

\begin{tcolorbox}[breakable,colback=teal!4,colframe=teal!35!black,title={Before you move on}]
What does \fa{تب} \fa{کردن} mean? (\textquotedblleft{}To run a fever\textquotedblright{}.) Say it once more.
\end{tcolorbox}
